\captionsetup{justification=raggedright,singlelinecheck=false}
\begin{longtable}{|p{0.5cm}|p{3cm}|p{2.3cm}|p{3.6cm}|p{3.6cm}|}
\caption{Studi Sebelumnya}
\label{tbl:studi-sebelumnya} \\
\hline
\textbf{No} &
\textbf{Judul Penelitian (Penulis)} &
\textbf{Metode} &
\textbf{Hasil Utama} &
\textbf{Keterbatasan} \\
\hline
\endfirsthead

\hline
\textbf{No} &
\textbf{Judul Penelitian (Penulis)} &
\textbf{Metode} &
\textbf{Hasil Utama} &
\textbf{Keterbatasan} \\
\hline
\endhead

\hline
\endfoot

1 &
\textit{Design and Implementation of an Android Spyware Application Using a Fake Game} \autocite{salih2020spyware} &
Perancangan sistem dan eksperimen &
Mengembangkan \textit{spyware} Android yang melakukan \textit{data crawling} terhadap kontak, pesan, \textit{log} panggilan, akun, dan lokasi, serta mengekstraksinya secara otomatis ke server \textit{backend} melalui REST API &
Berfokus pada Android saja, tidak membahas mekanisme \textit{stealth} tingkat lanjut atau deteksi oleh \textit{anti-spyware}. \\
\hline

2 &
\textit{Cyber Threat Intelligence through Automated Dark Web Crawling} \autocite{jegan2024aicrawler} &
Perancangan sistem dan eksperimen &
Mengimplementasikan \textit{web crawling} otomatis pada \textit{dark web} dengan ekstraksi konten, pemetaan tautan, dan deteksi kredensial yang terekspos menggunakan \textit{machine learning}. &
Fokus pada \textit{crawling} di \textit{dark web} untuk CTI, bukan \textit{crawling} internal pada perangkat korban seperti \textit{spyware}. \\
\hline

3 &
\textit{Stealthy Data Collection and Exfiltration in Pegasus Spyware} \autocite{kareem2024pegasus} &
Studi literatur teknis dan analisis arsitektur &
Mendeskripsikan mekanisme \textit{data collection} dan \textit{auto-exfiltration} \textit{spyware} Pegasus, mencakup akses pesan, email, media sosial, \textit{keylogging}, rekaman audio/video, lokasi \textit{real-time}, serta pengiriman terenkripsi ke C2. &
Tidak menyertakan implementasi uji empiris; sebagian informasi berasal dari laporan investigatif, bukan eksperimen langsung. \\
\hline

4 &
\textit{An Emerging Threat: Fileless Malware -- A Survey and Research Challenges} \autocite{sudhakar2020fileless} &
Studi literatur &
Menjelaskan teknik penyembunyian dan eksfiltrasi \textit{fileless}, termasuk penggunaan \textit{PowerShell}, \textit{registry abuse}, dan \textit{memory-resident payload}. &
Tidak membahas implementasi \textit{crawling} secara spesifik atau eksfiltrasi berbasis protokol tertentu. \\
\hline

\end{longtable}